% ═══════════════════════════════════════════════════════════════
% MUSTERLÖSUNG: Halbwertszeit (Physik Klasse 10)
% Identisches Layout wie AB, Lösungen in ROT
% ═══════════════════════════════════════════════════════════════

\documentclass[10pt, a4paper]{article}

\usepackage[utf8]{inputenc}
\usepackage[T1]{fontenc}
\usepackage[ngerman]{babel}
\usepackage{helvet}
\renewcommand{\familydefault}{\sfdefault}

\usepackage[margin=1.5cm, top=1.8cm, bottom=1.5cm]{geometry}
\usepackage{enumitem}
\usepackage{tabularx}
\usepackage{booktabs}
\usepackage{array}
\usepackage{multicol}

\usepackage{xcolor}
\usepackage{tcolorbox}
\usepackage{amssymb}
\usepackage{fancyhdr}
\usepackage{lastpage}

\setlist{nosep, leftmargin=1.5em}

% Farben
\definecolor{physik}{HTML}{2c5aa0}
\definecolor{grau}{HTML}{888888}
\definecolor{hellgrau}{HTML}{f0f0f0}
\definecolor{fehler}{HTML}{c62828}
\definecolor{fehlerbg}{HTML}{ffebee}
\definecolor{loesung}{HTML}{c62828}

\newcommand{\fachfarbe}{physik}

% Variablen
\newcommand{\thema}{Halbwertszeit}
\newcommand{\fach}{Physik}
\newcommand{\klasse}{Klasse 10}
\newcommand{\stunde}{Stunde 5-6}

% Kopf- und Fußzeile
\pagestyle{fancy}
\fancyhf{}
\renewcommand{\headrulewidth}{0.4pt}
\lhead{\small\fach{} \klasse{} -- \thema{} \textcolor{loesung}{\textbf{[MUSTERLÖSUNG]}}}
\rhead{\small Name: \rule{4cm}{0.4pt}}
\cfoot{\small\thepage/\pageref{LastPage}}

% Befehle
\newcommand{\lsg}[1]{\textcolor{loesung}{\textbf{#1}}}

\newtcolorbox{merksatz}[1][]{
    colback=hellgrau, colframe=\fachfarbe, boxrule=0.8pt, arc=0pt,
    left=6pt, right=6pt, top=6pt, bottom=6pt,
    fonttitle=\bfseries\small, title=#1
}

\newtcolorbox{fehlerbox}{
    colback=fehlerbg, colframe=fehler, boxrule=0.8pt, arc=0pt,
    left=6pt, right=6pt, top=4pt, bottom=4pt,
    fonttitle=\bfseries\small, title=Typische Fehler
}

\newcommand{\abschnitt}[1]{%
    \vspace{8pt}\noindent\textbf{\textcolor{\fachfarbe}{#1}}\vspace{4pt}
}

\begin{document}

\begin{center}
{\large\bfseries Arbeitsblatt: \thema{} -- \stunde{} \textcolor{loesung}{[MUSTERLÖSUNG]}}\\[2pt]
{\small\textcolor{grau}{Begleitend zum Lernpfad -- in den Hefter übernehmen}}
\end{center}
\vspace{-2pt}\hrule\vspace{8pt}

\begin{multicols}{2}
\small

% ─────────────────────────────────────────────────────────────
% KASTEN 1: Definition Halbwertszeit
% ─────────────────────────────────────────────────────────────

\abschnitt{Kasten 1: Definition Halbwertszeit}

\begin{merksatz}[Grundbegriff]
Die \textbf{Halbwertszeit} $T_{1/2}$ ist die Zeit, in der

\lsg{die Hälfte} der radioaktiven Kerne

zerfallen ist.

\vspace{6pt}

Nach einer Halbwertszeit sind noch \lsg{50} \%

der ursprünglichen Kerne vorhanden.

\vspace{6pt}

\textbf{Wichtig:} Die Halbwertszeit ist

\lsg{konstant / charakteristisch} für jedes Isotop.
\end{merksatz}

\vspace{6pt}

% ─────────────────────────────────────────────────────────────
% KASTEN 2: Das Zerfallsgesetz
% ─────────────────────────────────────────────────────────────

\abschnitt{Kasten 2: Das Zerfallsgesetz}

\begin{merksatz}[Formel]
\textbf{Zerfallsgesetz:}

\vspace{6pt}

$N(t) = N_0 \cdot \left(\frac{1}{2}\right)^{t/T_{1/2}}$

\vspace{6pt}

\renewcommand{\arraystretch}{1.4}
\begin{tabular}{|l|l|}
\hline
\textbf{Symbol} & \textbf{Bedeutung} \\
\hline
$N(t)$ & \lsg{Anzahl Kerne zur Zeit t} \\
\hline
$N_0$ & \lsg{Anfangsanzahl Kerne} \\
\hline
$t$ & \lsg{verstrichene Zeit} \\
\hline
$T_{1/2}$ & \lsg{Halbwertszeit} \\
\hline
\end{tabular}
\renewcommand{\arraystretch}{1}
\end{merksatz}

\vspace{6pt}

% ─────────────────────────────────────────────────────────────
% KASTEN 3: Beispiel-Halbwertszeiten
% ─────────────────────────────────────────────────────────────

\abschnitt{Kasten 3: Beispiel-Halbwertszeiten}

\begin{merksatz}[Übersicht]
Ergänze aus dem Lernpfad:

\vspace{4pt}
\renewcommand{\arraystretch}{1.4}
\begin{tabular}{|l|c|}
\hline
\textbf{Isotop} & \textbf{Halbwertszeit} \\
\hline
Radon-220 & \lsg{56 s} \\
\hline
Jod-131 & \lsg{8 Tage} \\
\hline
Cobalt-60 & \lsg{5,3 Jahre} \\
\hline
Cäsium-137 & \lsg{30 Jahre} \\
\hline
C-14 & \lsg{5730 Jahre} \\
\hline
Uran-238 & \lsg{4,5 Mrd. Jahre} \\
\hline
\end{tabular}
\renewcommand{\arraystretch}{1}
\end{merksatz}

\columnbreak

% ─────────────────────────────────────────────────────────────
% ÜBUNG 1: Zuordnung (AFB I)
% ─────────────────────────────────────────────────────────────

\abschnitt{Übung 1: Tabelle vervollständigen}

Fülle die Tabelle aus (Start: 1000 Kerne):

\vspace{4pt}
\renewcommand{\arraystretch}{1.4}
\begin{tabular}{|c|c|c|}
\hline
\textbf{Zeit} & \textbf{Anzahl Kerne} & \textbf{Anteil} \\
\hline
$0 \cdot T_{1/2}$ & 1000 & 100\% \\
\hline
$1 \cdot T_{1/2}$ & \lsg{500} & \lsg{50\%} \\
\hline
$2 \cdot T_{1/2}$ & \lsg{250} & \lsg{25\%} \\
\hline
$3 \cdot T_{1/2}$ & \lsg{125} & \lsg{12,5\%} \\
\hline
$4 \cdot T_{1/2}$ & \lsg{62,5} & \lsg{6,25\%} \\
\hline
\end{tabular}
\renewcommand{\arraystretch}{1}

\vspace{8pt}

% ─────────────────────────────────────────────────────────────
% ÜBUNG 2: Berechnungen (AFB II)
% ─────────────────────────────────────────────────────────────

\abschnitt{Übung 2: Berechnungen}

\textbf{a)} Ein Präparat enthält 8000 radioaktive Kerne.

Die Halbwertszeit beträgt 2 Stunden.

Wie viele Kerne sind nach 6 Stunden noch da?

\vspace{4pt}
\lsg{6h = 3 HWZ → 8000 → 4000 → 2000 → 1000 Kerne}

\vspace{6pt}

\textbf{b)} Jod-131 hat $T_{1/2} = 8$ Tage.

Von 1200 Kernen sind noch 150 übrig.

Wie viel Zeit ist vergangen?

\vspace{4pt}
\lsg{1200→600→300→150 = 3 HWZ → 3·8d = 24 Tage}

\vspace{8pt}

% ─────────────────────────────────────────────────────────────
% ÜBUNG 3: Fehler erklären (AFB III)
% ─────────────────────────────────────────────────────────────

\abschnitt{Übung 3: Fehler erklären (AFB III)}

\begin{fehlerbox}
Erkläre, warum diese Aussage falsch ist:

\vspace{4pt}

\textit{„Nach zwei Halbwertszeiten ist ein radioaktives Präparat komplett zerfallen."}

\vspace{4pt}
\lsg{Falsch! Nach 2 HWZ sind noch 25\% übrig. Nach 3 HWZ noch 12,5\%. Theoretisch erreicht die Kurve nie 0, praktisch ist die Aktivität irgendwann zu gering zum Messen.}
\end{fehlerbox}

\vspace{8pt}

% ─────────────────────────────────────────────────────────────
% ÜBUNG 4: Transfer (AFB III)
% ─────────────────────────────────────────────────────────────

\abschnitt{Übung 4: Anwendung C-14-Methode}

Ein Archäologe findet Holz mit 25\% des ursprünglichen C-14-Gehalts. Wie alt ist das Holz ungefähr?

($T_{1/2}$ von C-14 = 5730 Jahre)

\vspace{4pt}
\lsg{25\% = 1/4 = nach 2 HWZ}

\vspace{4pt}
\lsg{2 × 5730 Jahre = ca. 11.460 Jahre alt}

\end{multicols}

% ═══════════════════════════════════════════════════════════════
% SEITE 2: SIMULATION ZERFALLSKURVE
% ═══════════════════════════════════════════════════════════════

\newpage

\begin{center}
{\large\bfseries Simulationsprotokoll: Radioaktiver Zerfall \textcolor{loesung}{[MUSTERLÖSUNG]}}\\[2pt]
{\small\textcolor{grau}{Simulation: PhET Radioactive Dating Game}}
\end{center}
\vspace{-2pt}\hrule\vspace{12pt}

\abschnitt{Teil A: Zerfallskurve zeichnen}

\begin{merksatz}[Koordinatensystem]
Zeichne die Zerfallskurve für 1000 Ausgangskerne in das Koordinatensystem:

\vspace{8pt}
\lsg{Exponentiell abfallende Kurve: Punkte bei (0|1000), (1|500), (2|250), (3|125), (4|62,5), usw. Die Kurve nähert sich asymptotisch der x-Achse.}
\vspace{50pt}

Y-Achse: 0, 125, 250, 500, 750, 1000 \quad X-Achse: 0, 1, 2, 3, 4, 5, 6, 7

\end{merksatz}

\vspace{8pt}

\abschnitt{Teil B: Simulation beobachten}

Starte die PhET-Simulation und wähle verschiedene Isotope.

\vspace{6pt}

\begin{tabular}{|l|c|c|c|}
\hline
\textbf{Isotop} & \textbf{Halbwertszeit} & \textbf{Nach 2 HWZ (\%)} & \textbf{Nach 4 HWZ (\%)} \\
\hline
Isotop 1: \lsg{C-14} & \lsg{5730 J.} & \lsg{25\%} & \lsg{6,25\%} \\
\hline
Isotop 2: \lsg{U-238} & \lsg{4,5 Mrd. J.} & \lsg{25\%} & \lsg{6,25\%} \\
\hline
\end{tabular}

\vspace{12pt}

\abschnitt{Teil C: Auswertung}

\textbf{1. Was fällt dir bei allen Zerfallskurven auf (Form)?}

\vspace{4pt}
\lsg{Alle Kurven haben die gleiche Form: exponentieller Abfall. Sie werden nie zu Null, nähern sich aber asymptotisch der x-Achse.}

\vspace{8pt}

\textbf{2. Vergleiche ein kurzlebiges und ein langlebiges Isotop. Was ist gleich, was ist unterschiedlich?}

\vspace{4pt}
\lsg{Gleich: Die Form der Kurve (exponentiell), der prozentuale Anteil nach n Halbwertszeiten (50\%, 25\%, 12,5\%...)}

\vspace{4pt}
\lsg{Unterschiedlich: Die Zeitskala auf der x-Achse. Bei kurzlebigen Isotopen geht alles schneller.}

\vspace{8pt}

\textbf{3. Warum ist die Zerfallskurve keine Gerade, sondern eine Kurve?}

\vspace{4pt}
\lsg{Weil immer die HÄLFTE der noch vorhandenen Kerne zerfällt, nicht eine feste Anzahl. Je weniger Kerne noch da sind, desto weniger zerfallen pro Zeiteinheit.}

\vspace{4pt}
\lsg{Mathematisch: Es zerfällt ein konstanter ANTEIL, nicht eine konstante ANZAHL.}

\vspace{8pt}

\textbf{4. Erkläre: Warum kann man nie genau vorhersagen, WANN ein einzelner Kern zerfällt?}

\vspace{4pt}
\lsg{Der radioaktive Zerfall ist ein statistischer Prozess. Für einen einzelnen Kern ist es zufällig, wann er zerfällt.}

\vspace{4pt}
\lsg{Nur bei sehr vielen Kernen kann man Wahrscheinlichkeitsaussagen machen (Gesetz der großen Zahlen).}

\end{document}
